\begin{abstract}
Names in multi-agent simulations of language-model agents are usually treated as scaffolding.
We varied them as the experimental manipulation in a generational population of LLM agents
(Claude Haiku~5.5). The agents forage, hunt in pairs, form couples, endow their children, compete
with them for survival, and vote on their own scoring rule. Every arm saw identical information
and identical random draws. Only the handle each agent carried changed, and how that handle was
described. The evidence comes from four pre-registered studies: 480 analysed eight-generation
runs and about 367,000 model calls.

Any stable individual handle, whether a name or a bare ID and whoever gave it, raised pair
hunting two- to fivefold over no handle. Couple formation instead followed a parental frame. A
meaningless ID raised it when parents picked it, and equally when agents were merely told their
parents had given it. Only names composed by the parents raised investment in newborns, and
hunters with such names shared their kills most often. Arbitrary Blue/Green bands produced
in-group pairing and hunting. Yet a band added to an individual ID doubled pair hunting rather
than splitting agents into factions.

Set against human findings registered in advance, the identification and group-label effects
matched in direction. Two did not match. Namesaking did not raise parental investment, and agents
put about 3\% of their holdings into a public good, against the 40--60\% typical of people.

Exploratory analyses add three leads:
\begin{itemize}
  \item a world-assigned ID lowered couple formation below no handle at all;
  \item parents gave less to a child that the lot placed in the other band;
  \item a single sentence saying that children are named after a parent halved couple formation.
\end{itemize}

How agents are named, and how naming is described, shapes cooperation and family formation in
LLM agent populations.
\end{abstract}
